\chapter{Giới thiệu đề tài}

\section{Đặt vấn đề}

Trong những năm gần đây, trí tuệ nhân tạo (AI) đã trở thành một trong những công nghệ quan trọng thúc đẩy quá trình chuyển đổi số của doanh nghiệp. Các giải pháp AI được ứng dụng trong nhiều lĩnh vực như chăm sóc khách hàng, phân tích dữ liệu, tự động hóa quy trình nghiệp vụ và hỗ trợ ra quyết định.

Tuy nhiên, nhiều doanh nghiệp vẫn gặp khó khăn trong việc tiếp cận và triển khai các giải pháp AI do thiếu nguồn nhân lực chuyên môn, hạn chế về kinh nghiệm kỹ thuật và khó khăn trong việc tìm kiếm chuyên gia phù hợp. Trong khi đó, các chuyên gia AI và freelancer lại gặp khó khăn trong việc tiếp cận khách hàng có nhu cầu thực sự.

Từ thực tế đó, đề tài \textbf{AITasker – AI Marketplace Platform for AI Automation Services} được đề xuất nhằm xây dựng một nền tảng kết nối giữa doanh nghiệp và chuyên gia AI, hỗ trợ tìm kiếm, trao đổi và triển khai các dự án tự động hóa bằng AI.

\section{Lý do chọn đề tài}

Đề tài được lựa chọn dựa trên các lý do sau:

\begin{itemize}
    \item Nhu cầu ứng dụng AI trong doanh nghiệp ngày càng tăng.
    \item Thiếu các nền tảng chuyên biệt kết nối doanh nghiệp với chuyên gia AI.
    \item AI và tự động hóa đang là xu hướng công nghệ có tính ứng dụng cao.
    \item Đề tài cho phép áp dụng các kiến thức về phân tích, thiết kế và phát triển phần mềm.
\end{itemize}

\section{Mục tiêu đề tài}

Mục tiêu tổng quát của đề tài là xây dựng nền tảng AITasker nhằm kết nối doanh nghiệp có nhu cầu tự động hóa bằng AI với các chuyên gia AI.

Các mục tiêu cụ thể gồm:

\begin{itemize}
    \item Xây dựng hệ thống quản lý tài khoản và phân quyền người dùng.
    \item Hỗ trợ doanh nghiệp đăng tải yêu cầu dự án AI.
    \item Hỗ trợ chuyên gia AI tìm kiếm và nhận dự án phù hợp.
    \item Cung cấp chức năng gửi báo giá và trao đổi giữa các bên.
    \item Hỗ trợ theo dõi tiến độ thực hiện dự án.
    \item Tích hợp chức năng AI hỗ trợ gợi ý chuyên gia phù hợp.
\end{itemize}

\section{Đối tượng và phạm vi nghiên cứu}

\subsection{Đối tượng nghiên cứu}

Đề tài tập trung nghiên cứu các hệ thống kết nối dịch vụ trực tuyến, quy trình quản lý dự án AI và các kỹ thuật hỗ trợ ghép nối giữa doanh nghiệp và chuyên gia.

\subsection{Phạm vi nghiên cứu}

Đối tượng sử dụng hệ thống gồm:

\begin{itemize}
    \item Doanh nghiệp.
    \item Chuyên gia AI.
    \item Quản trị viên hệ thống.
\end{itemize}

Các chức năng chính bao gồm:

\begin{itemize}
    \item Đăng ký và đăng nhập.
    \item Quản lý hồ sơ người dùng.
    \item Đăng dự án AI.
    \item Gửi và quản lý báo giá.
    \item Chat trao đổi.
    \item Theo dõi tiến độ dự án.
    \item Đánh giá và phản hồi.
    \item Quản trị hệ thống.
\end{itemize}

\section{Cấu trúc báo cáo}

Nội dung báo cáo được tổ chức thành 6 chương như sau:

\begin{itemize}
    \item \textbf{Chương 1: Giới thiệu đề tài}

    Trình bày bối cảnh, lý do lựa chọn đề tài, mục tiêu nghiên cứu, phạm vi nghiên cứu, phương pháp thực hiện và ý nghĩa của đề tài.

    \item \textbf{Chương 2: Khảo sát và phân tích yêu cầu hệ thống}

    Phân tích bài toán, xác định các đối tượng sử dụng, khảo sát yêu cầu chức năng và phi chức năng, xây dựng các mô hình nghiệp vụ của hệ thống.

    \item \textbf{Chương 3: Phân tích và thiết kế hệ thống}

    Trình bày các sơ đồ UML như Use Case Diagram, Activity Diagram, Sequence Diagram, Class Diagram; thiết kế cơ sở dữ liệu và kiến trúc hệ thống.

    \item \textbf{Chương 4: Xây dựng hệ thống}

    Mô tả công nghệ sử dụng, triển khai cơ sở dữ liệu, xây dựng các chức năng chính và giao diện của hệ thống.

    \item \textbf{Chương 5: Kiểm thử và đánh giá hệ thống}

    Thực hiện kiểm thử chức năng, đánh giá kết quả đạt được, phân tích ưu điểm và hạn chế của hệ thống.

    \item \textbf{Chương 6: Kết luận và hướng phát triển}

    Tổng kết kết quả nghiên cứu, những đóng góp của đề tài và đề xuất các hướng phát triển trong tương lai.
\end{itemize}